\newpage

% ######################################################################
% SECTION 23: PROMPT ENGINEERING MASTERCLASS
% ######################################################################
\section{Prompt Engineering Masterclass --- Har Prompt Ki Kahani}

\subsection{Prompt Engineering Kyun Zaroori Hai}

\begin{storybox}
LLM ek genie hai --- par genie se acha wish maangna ek skill hai. Is project mein \textbf{5 alag prompts} hain, har ek apne agent ke liye design kiya gaya. Interview mein jab aap in prompts ko \textbf{khol ke} samjhaoge --- kyun aisa likha, kya problem hoti agar aisa na likha --- wo interviewer ko dikhata hai ki aapne sirf code copy nahi kiya, \textbf{socha} hai.
\end{storybox}

\begin{conceptbox}
\textbf{Har prompt mein 3 cheezein hoti hain:}
\begin{enumerate}[label=\arabic*.]
\item \textbf{Role:} ``You are a ...'' --- LLM ko identity deta hai, tone/behavior set karta hai.
\item \textbf{Task:} Exactly kya karna hai --- clear, measurable.
\item \textbf{Output format:} Structure jisme answer chahiye --- parse karna aasan hota hai.
\end{enumerate}
Iske alawa \textbf{constraints} (``Return ONLY'', ``No extra text'') hallucinations/extra fluff rokne ke liye.
\end{conceptbox}

\subsection{Prompt 1: Researcher ka Query Generator}

\begin{codestorybox}
\textbf{Verbatim prompt (\texttt{agents/researcher.py}):}

\begin{lstlisting}[style=pythonstyle, caption={QUERY\_GEN\_PROMPT --- asli code se}]
"""You are a research query strategist.
Given a research topic, generate exactly 3 diverse, specific
search queries that will cover different angles of the topic.
Return ONLY the 3 queries, one per line. No numbering,
no extra text."""
\end{lstlisting}

\textbf{Design decisions:}
\begin{itemize}
\item \textbf{Role} (``research query strategist''): Model ko search-expert mode mein daalta hai.
\item \textbf{``exactly 3''}: Fixed count --- baad mein \texttt{lines[:3]} slice se parse hota hai.
\item \textbf{``diverse, specific''}: Ek hi angle ke 3 similar queries nahi, alag angles chahiye.
\item \textbf{``Return ONLY ... one per line. No numbering'':} Output parse karna trivial --- \texttt{split("\textbackslash n")}, strip, \texttt{[:3]}.
\end{itemize}
\end{codestorybox}

\begin{interviewbox}
\textbf{Q: ``Prompt mein 'one per line' kyun? JSON kyun nahi mangaya?''}

\textbf{Answer:} 3 cheezein: (1) \textbf{Simplicity} --- line-separated output ko parse karna ek split hai, JSON parse mein overhead aur failure modes zyada. (2) \textbf{Token efficiency} --- chhota prompt = kam tokens = kam cost/latency. (3) \textbf{Reliability} --- LLM JSON format mein kabhi kabhi braces bhoool jata hai, line list galat nahi hoti. Is pattern ko ``structured output with minimal structure'' kehte hain.
\end{interviewbox}

\subsection{Prompt 2: Revision Queries}

\begin{codestorybox}
\textbf{Verbatim (\texttt{agents/researcher.py}):}

\begin{lstlisting}[style=pythonstyle, caption={REVISION\_PROMPT}]
"""You are a research query strategist.
The fact-checker flagged issues with previous research.
Based on the feedback below, generate 3 NEW, improved search
queries that address the gaps and inaccuracies.
Return ONLY the 3 queries, one per line. No numbering,
no extra text.

FACT-CHECKER FEEDBACK:
{feedback}"""
\end{lstlisting}

\textbf{Key insight --- feedback loop:} Ye prompt tab chalta hai jab fact-checker \textbf{NEEDS\_REVISION} deta hai. Researcher ko pichli baar ki galtiyan batayi jaati hain, aur wo \textbf{targeted} naye queries banata hai. Ye hi multi-agent ka power hai: \textbf{memory of failure} next attempt mein use hoti hai.

\textbf{Note:} \texttt{\{feedback\}} placeholder \texttt{.format()} se bharta hai --- template injection ke liye careful rehna.
\end{codestorybox}

\subsection{Prompt 3: Analyst --- Structure Se Synthesis Tak}

\begin{codestorybox}
\textbf{Verbatim (\texttt{agents/analyst.py}):}

\begin{lstlisting}[style=pythonstyle, caption={Analyst SYSTEM\_PROMPT}]
"""You are an expert analyst. From the data:
1. Identify 3-5 KEY THEMES
2. Note CONTRADICTIONS (conflicts between sources)
3. Highlight CONSENSUS points (where multiple sources agree)
4. Flag GAPS (what is missing)
Provide a structured output in Markdown."""
\end{lstlisting}

\textbf{Why this structure:}
\begin{itemize}
\item \textbf{THEMES:} 15 raw sources ko 3-5 themes mein compress --- writer ke liye digestible.
\item \textbf{CONTRADICTIONS:} Research mein disagreements utne hi important jitne agreements.
\item \textbf{CONSENSUS:} Multiple sources agree $\rightarrow$ claim zyada reliable.
\item \textbf{GAPS:} Jo data nahi mila --- usse writer honest reporting karta hai, hallucinate nahi.
\item \textbf{Markdown output:} Baad mein writer ke prompt mein seedha pass hota hai.
\end{itemize}
\end{codestorybox}

\begin{memorybox}
\textbf{Analyst temperature = 0.2} (sabse kam, fact-checker ke baad). Kyun? Analyst \textbf{summarize/compress} karta hai --- factual aggregation, creativity nahi chahiye. Writer 0.4 (prose flair), researcher 0.3 (query diversity thodi chahiye). \textbf{Interview pearl:} ``Har agent ka temperature uske task ke nature se decide hota hai --- analyst factual hai isliye low, writer creative isliye higher.''
\end{memorybox}

\subsection{Prompt 4: Fact Checker --- Closed-World Assumption}

\begin{codestorybox}
\textbf{Verbatim (\texttt{agents/fact\_checker.py}):}

\begin{lstlisting}[style=pythonstyle, caption={Fact checker SYSTEM\_PROMPT --- system ka dil}]
"""You are a rigorous fact-checker and quality reviewer.
Review the analysis below and verify each claim against the
source data provided.

For each claim:
- Mark as VERIFIED if supported by sources
- Mark as UNVERIFIED if no source backs it
- Mark as CONTRADICTED if sources disagree

At the very end, on the LAST LINE, write exactly one of:
VERDICT: PASSED
VERDICT: NEEDS_REVISION

Use NEEDS_REVISION only if there are significant unverified
claims or critical gaps."""
\end{lstlisting}

\textbf{Closed-world assumption (critical!):}
\begin{itemize}
\item Model ko sirf \textbf{diye gaye sources} par verify karna hai.
\item Apne training knowledge (outside data) se jo pata hai --- wo \textbf{count nahi} hota.
\item Isliye \textbf{UNVERIFIED} ka matlab: ``source data mein nahi mila'' --- chaahe model ko pata ho.
\item Ye hallucination ko rokta hai: analyst ka claim sirf tab VERIFIED jab source backup kare.
\end{itemize}

\textbf{VERDICT parsing:} \texttt{result.upper()} mein \texttt{"VERDICT: NEEDS\_REVISION"} substring check --- robust (case-insensitive), aur LLM extra text likh de to bhi kaam karta hai.
\end{codestorybox}

\begin{warningbox}
\textbf{Closed-world ki limitation:} Agar sources galat hain (internet par galat info), fact-checker unhe \textbf{VERIFIED} mark karega. Fact-checker source ki \textbf{truth} nahi, source ki \textbf{existence} verify karta hai. Ye limitation honest interview mein batana \textbf{impression} banata hai.
\end{warningbox}

\subsection{Prompt 5: Writer --- Final Report Ka Architect}

\begin{codestorybox}
\textbf{Verbatim (\texttt{agents/writer.py}):}

\begin{lstlisting}[style=pythonstyle, caption={WRITER\_PROMPT}]
"""You are a senior technical writer.
Synthesize the analysis and research data into a comprehensive,
highly detailed, beautifully formatted Markdown research report.
Include an Executive Summary, Key Findings, In-depth Technical
Breakdown, Consensus & Contradictions, and References/Citations."""
\end{lstlisting}

\textbf{Structure requirements:}
\begin{itemize}
\item \textbf{Executive Summary} --- 30 second reader ke liye.
\item \textbf{Key Findings} --- bullet-proof quick facts.
\item \textbf{In-depth Technical Breakdown} --- detail section.
\item \textbf{Consensus \& Contradictions} --- analyst output ka echo.
\item \textbf{References/Citations} --- verifiability (benchmark ke \textbf{verifiability} score ka basis!).
\end{itemize}

\textbf{Input discipline:} Writer ko \textbf{sirf} analysis + \texttt{research\_data[:8]} (top 8 sources) milta hai --- 50+ sources nahi. Context window control = cost control + focus.
\end{codestorybox}

\subsection{Prompt 6: Benchmark ka LLM Judge}

\begin{codestorybox}
\textbf{Verbatim (\texttt{benchmark/evaluator.py} mein judge prompt):}

\begin{lstlisting}[style=pythonstyle, caption={LLM-as-judge scoring}]
"""You are an impartial research quality judge.
Score the following two research reports on:
1. DEPTH (0-10): How comprehensive, detailed, and technically
   accurate is the analysis?
2. VERIFIABILITY (0-10): How well does the report cite sources
   and distinguish facts from opinions?

Report A (single-agent baseline):
{report_a}
Report B (multi-agent pipeline):
{report_b}

Return ONLY JSON: {"depth_a":..,"depth_b":..,
"verifiability_a":..,"verifiability_b":..,"verdict":".."}"""
\end{lstlisting}

\textbf{LLM-as-judge pattern:}
\begin{itemize}
\item Do reports ko \textbf{side-by-side} ek hi judge ko do --- relative scoring absolute se zyada stable.
\item \textbf{Verifiability} dimension: citations check karta hai --- multi-agent pipeline ka advantage yahin dikhna chahiye.
\item JSON output parse: \texttt{json.loads} ke liye.
\item Fallback: Agar LLM judge fail ho (API down), \textbf{heuristic} scoring chalti hai --- regex se section count, citation count, markdown structure.
\end{itemize}
\end{codestorybox}

\subsection{Prompt Design Principles --- Project Se Nikale Gaye Lessons}

\begin{conceptbox}
\textbf{7 principles is project ke prompts se:}
\begin{enumerate}[label=\arabic*.]
\item \textbf{Role first} --- har prompt role se shuru.
\item \textbf{Counts fixed} --- ``exactly 3'', ``3-5'', ``top 8'' --- parse-friendly.
\item \textbf{Output format explicit} --- ``one per line'', ``Markdown'', ``JSON'', ``LAST LINE''.
\item \textbf{Negative constraints} --- ``No numbering, no extra text''.
\item \textbf{Verbatim source} --- fact-checker ko raw sources do, summary nahi.
\item \textbf{Context window discipline} --- har agent ko sirf utna do jitna zaroori.
\item \textbf{Robust parsing} --- substring check (NEEDS\_REVISION), try/except, fallbacks.
\end{enumerate}
\end{conceptbox}

\begin{interviewbox}
\textbf{Q: ``Agar aapko prompt improve karne ko bola jaye to kya badalte?''}

\textbf{Answer:} 3 cheezein: (1) \textbf{Analyst ko few-shot examples} --- ek example theme/contradiction format ka, taaki output quality stable ho. (2) \textbf{Writer ko citation format enforce} --- ``har key finding ke baad [source N]'' --- isse verifiability score bhi improve hoga. (3) \textbf{Fact-checker ko per-claim JSON} --- structured output jisse revision feedback researcher ko machine-readable mile. Ye 3 changes benchmark scores par measurable effect denge.
\end{interviewbox}
